\documentclass[14pt]{extarticle} % увеличенный базовый шрифт
\usepackage{amsmath}
\usepackage{mathtext}
\usepackage{array}
\usepackage{fancyvrb}
\usepackage{minted}
\usepackage{graphicx}
\usepackage[hidelinks]{hyperref}
\usepackage[utf8]{inputenc}
\usepackage[T2A]{fontenc}
\usepackage[russian]{babel}
\usepackage{ragged2e} % для выравнивания текста



% Выравнивание текста по ширине
\justifying

% Увеличение межстрочного интервала для лучшей читаемости
\usepackage{setspace}
\singlespacing

% Настройка полей (опционально, для лучшего форматирования)
\usepackage{geometry}
\geometry{a4paper, left=2.5cm, right=2.5cm, top=2cm, bottom=2cm}

\setlength{\parindent}{0pt}
\setlength{\parskip}{\baselineskip}



\begin{document}

\begin{titlepage}
    \centering % Центрирование всего содержимого
    
    \vspace*{\fill} % Заполнение вертикального пространства
    
    {\LARGE\bfseries Использование генетических алгоритмов и
нейронных сетей в анализе деформаций стопы\par} % Название предмета
    \vspace{2cm}
    {\Large \par} % Подзаголовок
    \vspace{0.5cm}
    {\Large 2 курс \par} % Подзаголовок
    \vspace{0.5cm}
    {\Large выполнил Кочергин Антон\par} % Подзаголовок
    \vspace{1cm}
    {\large Научный руководитель:\par}
    {\Large\bfseries Пантелеев \par} % Имя лектора
   
    
    \vspace{4cm}
    {\large 2025/2026 учебный год\par} % Год
    
    \vspace*{\fill} % Заполнение вертикального пространства
\end{titlepage}

% Сбрасываем центрирование для основного текста
\raggedright % Выравнивание по левому краю
\justifying  % Выравнивание по ширине (если нужно)

\newpage

\tableofcontents

\newpage

\setlength{\parindent}{0pt}


\section{Лекция №1 }



\end{document}